\subsection{零点定理}

命题 1. 设 A 为域 \(k\) 上的一个有限生成代数，\(L\) 为 \(k\) 的代数闭包。

\begin{enumerate}
\def\labelenumi{(\roman{enumi})}
\item
  若 \(\mathbf{A} \neq \{0\}\)，则存在一个 \(k\)-同态（从 \(\mathbf{A}\) 到 \(\mathbf{L}\)）。
\item
  设 \(f_{1}, f_{2}\) 为两个 \(k\)-同态（从 \(\mathbf{A}\) 到 \(\mathbf{L}\)）。要使 \(f_{1}\) 与 \(f_{2}\) 有相同的核，必须且只需存在一个 \(k\)-自同构 \(s\)（\(\mathbf{L}\) 的），使 \(f_{2} = s \circ f_{1}\)。
\item
  设 \(\mathfrak{a}\) 为 \(\mathbf{A}\) 的一个理想。要使 \(\mathfrak{a}\) 是极大的，必须且只需它是一个 \(k\)-同态（从 \(\mathbf{A}\) 到 \(\mathbf{L}\)）的核。
\item
  要使元素 \(x\)（\(\mathbf{A}\) 的）满足 \(f(x) = 0\)（对每个 \(k\)-同态 \(f\)，从 \(\mathbf{A}\) 到 \(\mathbf{L}\)），必须且只需 \(x\) 是幂零的。
\end{enumerate}

断言 (i) 由第 1 节定理 1 的推论 3 得出，应用时把 A 换成 k，B 换成 A，b 换成 B 的单位元，f 换成 k 到 L 的典范单射。

若 \(f\) 是一个 \(k\)-同态（从 \(\mathbf{A}\) 到 \(\mathbf{L}\)），则 \(f(\mathbf{A})\) 是 \(\mathbf{L}\) 的一个包含 \(k\) 的子环；由于 \(\mathbf{L}\) 是 \(k\) 的代数扩张，\(f(\mathbf{A})\) 是一个域（《代数》第 V 章 § 3 第 2 节命题 3），而若 \(\mathfrak{a}\) 是 \(f\) 的核，则 \(\mathbf{A} / \mathfrak{a}\)（同构于 \(f(\mathbf{A})\)）因此是一个域，这就证明 \(\mathfrak{a}\) 是极大的。反之，若 \(\mathfrak{a}\) 是 \(\mathbf{A}\) 的一个极大理想，则由 (i) 可知存在一个 \(k\)-同态（从 \(\mathbf{A} / \mathfrak{a}\) 到 \(\mathbf{L}\)），从而存在一个 \(k\)-同态（\(\mathbf{A}\) 到 \(\mathbf{L}\)），其核 \(\mathfrak{b}\) 包含 \(\mathfrak{a}\)；但由于 \(\mathfrak{a}\) 极大，\(\mathfrak{b} = \mathfrak{a}\)；这就证明了 (iii)。

现在证明 (ii)。若 \(s\) 是一个 \(k\)-自同构（\(L\) 的），满足 \(f_2 = s \circ f_1\)，则显然 \(f_1\) 与 \(f_2\) 有相同的核。反之，设 \(f_1\) 与 \(f_2\) 有相同的核；那么存在一个 \(k\)-同构 \(s_0\)（从域 \(f_1(A)\) 到域 \(f_2(A)\) 上），使 \(f_2 = s_0 \circ f_1\)；但由《代数》第 V 章 § 6 第 3 节命题 7，\(s_0\) 可延拓为一个 \(k\)-自同构 \(s\)（\(L\) 的），从而 \(f_2 = s \circ f_1\)。

最后，若 \(x \in \mathbf{A}\) 满足 \(x^n = 0\)，则对每个 \(k\)-同态 \(f\)（从 \(\mathbf{A}\) 到 \(\mathbf{L}\)），有 \((f(x))^n = f(x^n) = 0\)，从而 \(f(x) = 0\)，因为 \(\mathbf{L}\) 是域。反之，设 \(x \in \mathbf{A}\) 不是幂零的；那么 \(\mathbf{A}[x^{-1}]\) 是一个不退化为 0 的有限生成 A-代数（从而是有限生成 \(k\)-代数）（第 II 章 §2 第 1 节注 3），于是由 (i) 存在一个 \(k\)-同态 \(g\)（从 \(\mathbf{A}[x^{-1}]\) 到 \(\mathbf{L}\)）。若 \(j: \mathbf{A} \to \mathbf{A}[x^{-1}]\) 是典范同态，则 \(f = g \circ j\) 是一个 \(k\)-同态（从 \(\mathbf{A}\) 到 \(\mathbf{L}\)），且 \(f(x)g(1/x) = g(x/1)g(1/x) = g(1) = 1\)，故 \(f(x) \neq 0\)。

设 \(k\) 为域，\(\mathbf{L}\) 为 \(k\) 的一个扩域；元素 \(\mathbf{x} = (x_1, \ldots, x_n)\)（\(\mathbf{L}^n\) 中的）称为 \(\mathbf{L}^n\) 中的一个零点，即理想 \(\mathfrak{r}\)（多项式环 \(k[\mathbf{X}_1, \ldots, \mathbf{X}_n]\) 的理想）的零点，如果

\[
\mathrm{P} (\mathbf {x}) = \mathrm{P} (x _ {1}, \dots , x _ {n}) = 0
\]

对所有 \(\mathbf{P} \in \mathfrak{r}\)。

引理 1. 设 A 为域 \(k\) 上的一个有限生成代数，\((a_i)_{1 \leqslant i \leqslant n}\) 为它的一个生成系，\(\mathfrak{r}\) 为诸 \(a_i\) 之间的、系数在 \(k\) 中的代数关系理想（《代数》第 IV 章 §2 第 1 节）。对每个扩域 \(\mathbf{L}\)（\(k\) 的），映射 \(f \mapsto (f(a_i))_{1 \leqslant i \leqslant n}\) 是从 A 到 L 的 \(k\)-同态所成之集到 \(\mathfrak{r}\) 在 \(L^n\) 中的零点所成之集上的一个双射。

存在唯一的一个 \(k\)-代数同态 \(h\)，从 \(k[\mathbf{X}_1, \ldots, \mathbf{X}_n]\) 到 \(\mathbf{A}\) 上，使 \(h(\mathbf{X}_i) = a_i\)（对 \(1 \leqslant i \leqslant n\)），且按定义 \(\mathfrak{r}\) 是 \(h\) 的核。映射 \(f \mapsto f \circ h\) 是诸 \(k\)-同态（从 \(A\) 到 \(L\)）所成之集到诸 \(k\)-同态（从 \(k[X_1, \ldots, X_n]\) 到 \(L\)、在 \(r\) 上为零）所成之集上的一个双射。对每个多项式 \(P \in k[X_1, \ldots, X_n]\) 与每个元素 \(\mathbf{x} = (x_1, \ldots, x_n) \in L^n\)，记 \(h_{\mathbf{x}}(P) = P(\mathbf{x})\)；于是映射 \(\mathbf{x} \mapsto h_{\mathbf{x}}\) 是从 \(L^n\) 到诸 \(k\)-同态（从 \(k[X_1, \ldots, X_n]\) 到 \(L\)）所成之集上的一个双射（这样的同态由它在诸 \(X_i\)（\(1 \leqslant i \leqslant n\)）处的值确定）；说 \(h_{\mathbf{x}}\) 在 \(r\) 上为零，就是说 \(\mathbf{x}\) 是 \(r\) 在 \(L^n\) 中的一个零点，由此即得引理。

若把命题 1 应用于代数 \(\mathbf{A} = k[\mathbf{X}_1, \ldots, \mathbf{X}_n] / \mathfrak{r}\)，其中 \(\mathfrak{r}\) 是 \(k[\mathbf{X}_1, \ldots, \mathbf{X}_n]\) 的一个异于整个环的理想，则由引理 1 得到下述命题：

命题 2（Hilbert 零点定理）. 设 \(k\) 为域，\(L\) 为 \(k\) 的一个代数闭包。

\begin{enumerate}
\def\labelenumi{(\roman{enumi})}
\item
  每个理想 \(\mathfrak{r}\)（\(k[\mathbf{X}_1, \ldots, \mathbf{X}_n]\) 的），若不含 1，则在 \(\mathbf{L}^n\) 中至少有一个零点。
\item
  设 \(\mathbf{x} = (x_1, \ldots, x_n)\)、\(\mathbf{y} = (y_1, \ldots, y_n)\) 为 \(\mathbf{L}^n\) 的两个元素；要使 \(k[\mathbf{X}_1, \ldots, \mathbf{X}_n]\) 中在 \(\mathbf{x}\) 处为零的多项式所成之集与 \(k[\mathbf{X}_1, \ldots, \mathbf{X}_n]\) 中在 \(\mathbf{y}\) 处为零的多项式所成之集相同，必须且只需存在一个 \(k\)-自同构 \(s\)（\(\mathbf{L}\) 的），使 \(y_i = s(x_i)\)（对 \(1 \leqslant i \leqslant n\)）。
\item
  要使理想 \(\mathfrak{a}\)（\(k[\mathbf{X}_1, \ldots, \mathbf{X}_n]\) 的）是极大的，必须且只需存在一个 \(\mathbf{x}\)（在 \(\mathbf{L}^n\) 中），使 \(\mathfrak{a}\) 是 \(k[\mathbf{X}_1, \ldots, \mathbf{X}_n]\) 中在 \(\mathbf{x}\) 处为零的多项式所成之集。
\item
  要使多项式 \(\mathbf{Q}\)（\(k[\mathbf{X}_1, \ldots, \mathbf{X}_n]\) 的）在 \(\mathbf{L}^n of\) 中、理想 \(\mathfrak{r}\)（\(k[\mathbf{X}_1, \ldots, \mathbf{X}_n]\) 的）的零点所成之集上为零，必须且只需存在一个整数 \(m > 0\) 使 \(\mathbf{Q}^m \in \mathfrak{r}\)。
\end{enumerate}

